\documentclass[11pt]{article}
\usepackage[a4paper,margin=26mm]{geometry}
\usepackage{amsmath,amssymb,amsthm}
\usepackage[T1]{fontenc}
\usepackage{booktabs}
\usepackage{xurl}
\usepackage{hyperref}
\hypersetup{colorlinks=true,linkcolor=blue,urlcolor=blue,citecolor=blue,
 pdftitle={Proofs of eleven congruence conjectures of Hanna from the OEIS, and the resolution of three more},
 pdfauthor={Horace Dong}}
\newtheorem{theorem}{Theorem}
\newtheorem{lemma}[theorem]{Lemma}
\theoremstyle{remark}
\newtheorem{remark}[theorem]{Remark}
\newcommand{\F}{\mathbb F}
\newcommand{\Z}{\mathbb Z}
\newcommand{\N}{\mathbb N}
\newcommand{\ord}{\operatorname{ord}}
\title{Proofs of eleven congruence conjectures of Hanna\\
from the OEIS, and the resolution of three more}
\author{Horace Dong\\ \small Independent researcher\\
\small\texttt{maxwelldhx@gmail.com}\\
\small ORCID: \href{https://orcid.org/0009-0002-7554-7548}{0009-0002-7554-7548}}
\date{September 2026}
\begin{document}
\maketitle
\begin{abstract}
Many entries of the On-Line Encyclopedia of Integer Sequences contributed by
Paul D.~Hanna define a power series by a functional equation and conjecture a
congruence for its coefficients. We prove eleven such conjectures:
\begin{itemize}
\item the modulo-three conjecture of A389472;
\item the parity conjectures of A120566, A240998, A273958, A295762, A301933,
 A338633, A338634 and A388734;
\item the modulo-three conjectures of A274479 and A377100.
\end{itemize}
We also resolve three more:
\begin{itemize}
\item for A184894 and A107099 we prove the conjectured vanishing modulo~$3$
 outside the exceptional indices. We also show that the stronger reading,
 non-vanishing at every exceptional index, is false. It fails at the
 OEIS index $365$ of A184894, and at the coefficient of $x^{729}$ (OEIS
 index $364$) of A107099;
\item for A361047 we prove the intended statement, which the entry states
 with an index shift.
\end{itemize}
The proofs reduce the functional equations modulo a prime and use the
Frobenius identity, formal derivatives, explicit solutions with an $x$-adic
uniqueness argument, and Lucas's theorem. Every theorem is stated for every
integer power series satisfying the defining equation, and every theorem is
formally verified in Lean~4. The proofs were found by AI coding agents
working under the author's direction; their role is described at the end of
the paper.
\end{abstract}

\section{Introduction}
Throughout, $A(x)=\sum_{n\ge0}a(n)x^n\in\Z[[x]]$, $A'$ is its formal
derivative, and for a series $h$ with $h(0)=0$ we write $A(h)$ or $A\circ h$
for the composition $\sum_n a(n)h^n$. The results are summarised in
Table~\ref{tab:summary}; each conjecture is stated in the corresponding
entry~\cite{oeis}. The theorems hold for every integer solution of the
defining equation, together with the normalisation of the first terms fixed
by the entry where the equation alone does not determine it. The equation and
these initial terms determine the sequence by a triangular recursion with
integer solutions. There, the new coefficient appears with coefficient $1$
or~$2$. Where it appears with~$2$, the required divisibility holds: for A274479
by Lemma~\ref{lem:frob} modulo~$2$, and for A107099 by writing $A=x+2B$. For
the continued fractions, the tails exist as $x$-adic limits. Hence the
theorems apply to the sequences in the OEIS. These existence arguments are
not formalised; the Lean theorems hold for every integer solution.

\begin{table}[ht]\centering\small
\begin{tabular}{@{}lll@{}}\toprule
Entry & Defining equation & Result\\\midrule
A389472 & $A(x^2+x^3)=x^2(1+A)$ & $3\mid a(3n-1)$, $n>1$\\
A240998 & $A^2=x+A(x+2x^2)$ & $a(n)$ odd $\iff n=2^k$ ($n\ge1$)\\
A295762 & $A(x-2A(x^2))=x+A(x^2)$ & $a(n)$ odd $\iff n=2^k$ ($n\ge1$)\\
A273958 & $xA+x^2A^2=C^2$, $C=x+C^2$ & $a(n)$ odd $\iff n=2\cdot4^k-1$\\
A301933 & $A=x(1+4AA')/(1+AA')$ & $a(n)$ odd $\iff n=2^k$\\
A338633 & $1=A-x/(A-2^3x/(A-3^3x/\cdots))$ & $a(n)$ odd $\iff n=2^k$ ($n\ge1$)\\
A338634 & $1=A-x/(A-2^4x/(A-3^4x/\cdots))$ & $a(n)$ odd $\iff n=2^k$ ($n\ge1$)\\
A388734 & $A=1+xA^2+x^2(1-x)A^3$ & all $a(n)$ odd\\
A120566 & $A=A\circ A-x\,A\circ A\circ A$ & all $a(n)$ odd ($n\ge1$)\\
A377100 & $A=A(x^3)/A(x^2)+A^2$ & $a(n)\equiv1\pmod3$ ($n\ge1$)\\
A274479 & $A^2=A\bigl(x^2/(1-2x-4x^2)\bigr)$ & $a(n)\equiv1\pmod3$ ($n\ge1$)\\\midrule
A361047 & $A(x-x^3A'^2)=x$ & $[x^m]A\equiv[m=3^k]\pmod3$\\
A184894 & iterates of $x+x^3$ & vanishing proved; non-vanishing fails at index 365\\
A107099 & $A\circ A=x+4x^3$ & vanishing proved; non-vanishing fails at $x^{729}$\\\bottomrule
\end{tabular}
\caption{Summary. Normalisations and precise statements are given in the
theorems below.}\label{tab:summary}
\end{table}

\paragraph{Prior work.} The modulo-two conjecture of A389472 was proved
earlier by Perez Fontelles~\cite{pf389472} and, independently, in the
\texttt{trureturing} repository~\cite{tru389472}; both leave the modulo-three
conjecture open. We found no earlier proof of any of the results in this
paper in the OEIS entries and their histories, in the literature, or in the
public repositories of proofs of OEIS conjectures that we searched (among
them~\cite{pfrepo,trurepo}). This reports the coverage of our searches, not a
guarantee of priority.

\section{Tools}
Reduction of the coefficients modulo a prime $p$ is a ring homomorphism
$\Z[[x]]\to\F_p[[x]]$; it commutes with formal differentiation and with
composition by series with zero constant term. We write $F=\overline A$ for
the reduction and $f(n)$ for its coefficients.

\begin{lemma}[Frobenius]\label{lem:frob}
For $F\in\F_p[[x]]$ we have $F(x)^p=F(x^p)$.
\end{lemma}
\begin{proof}
The map $G\mapsto G^p$ is a continuous ring endomorphism of $\F_p[[x]]$
that fixes $\F_p$ and sends $x$ to $x^p$.
\end{proof}

\begin{lemma}\label{lem:fixed}
Let $K$ be a field, let $h\in xK[[x]]$ have order at least two, and let
$H\in K[[x]]$ satisfy $H(h)=H$. Then $H$ is constant.
\end{lemma}
\begin{proof}
Put $B=H-H(0)$, so $B(h)=B$. If $B\ne0$, its order $d$ is positive, while
$\ord B(h)\ge2d>d$.
\end{proof}

\begin{lemma}[dyadic recursion]\label{lem:dyadic}
Let $F\in\F_2[[x]]$ satisfy $F(0)=0$ and $F=x+F^2$. Then for $n\ge1$,
$f(n)=1$ if and only if $n$ is a power of~$2$.
\end{lemma}
\begin{proof}
By Lemma~\ref{lem:frob}, $F=x+F(x^2)$, so $f(n)=[n=1]+[2\mid n]f(n/2)$ for
$n\ge1$. By strong induction, $f(1)=1$, $f(n)=0$ for odd $n\ge3$, and
$f(2m)=f(m)$.
\end{proof}

We also use the following observation repeatedly: if $D\in\F_p[[x]]$,
$x^{n_0}\mid D$, and $x^n\mid D$ implies $x^{n+1}\mid D$ for all $n\ge n_0$,
then $D=0$.

\section{The modulo-three conjecture of A389472}
\begin{theorem}\label{thm:389472}
If $A(x^2+x^3)=x^2(1+A(x))$, then $a(3n-1)\equiv0\pmod3$ for every $n>1$.
\end{theorem}
\begin{proof}
Let $F=\overline A\in\F_3[[x]]$ and $h=x^2+x^3$. Then $F(h)=x^2(1+F)$ and
$h'=2x+3x^2=-x$. Differentiating, $-xF'(h)=-x(1+F)+x^2F'$, so
$F'(h)=1+F-xF'$ after cancelling $-x$. Differentiating again, the terms $F'$
cancel and $-xF''(h)=-xF''$, so $F''(h)=F''$. By Lemma~\ref{lem:fixed},
$F''$ is constant. For $n>1$ the coefficient of $x^{3n-3}$ in $F''$ is
therefore zero, and it equals $(3n-1)(3n-2)f(3n-1)=2f(3n-1)$.
\end{proof}

\section{Parity via the dyadic recursion}
\begin{theorem}\label{thm:parity}
In each of the following cases, $a(n)$ is odd for $n\ge1$ if and only if $n$
is a power of~$2$:
\begin{enumerate}
\item (A240998) $A^2=x+A(x+2x^2)$;
\item (A295762) $a(0)=0$ and $A(x-2A(x^2))=x+A(x^2)$;
\item (A301933) $a(0)=0$ and $A(1+AA')=x(1+4AA')$;
\item (A338633, A338634) $A$ satisfies the continued fraction
 $1=A-x/(A-2^ex/(A-3^ex/(A-\cdots)))$ with $e=3$, resp.\ $e=4$, in the
 following sense: there are series $T_0,T_1,\dots$ with constant terms
 $\pm1$ such that $(A-1)T_0=x$ and $(A-T_k)T_{k+1}=(k+2)^ex$ for all
 $k\ge0$.
\end{enumerate}
Moreover (A273958), if $C=x+C^2$ and $xA+x^2A^2=C^2$, then $a(n)$ is odd if
and only if $n=2\cdot4^k-1$ for some $k\ge0$.
\end{theorem}
\begin{proof}
(1) Modulo~$2$, $(x+2x^2)^k\equiv x^k$, so $F^2=x+F$. Then $G=F-f(0)$
satisfies $G(0)=0$ and $G^2=F^2-f(0)=x+G$, and Lemma~\ref{lem:dyadic}
applies to $G$.
(2) Modulo~$2$ the inner series reduces to $x$, so $F=x+F(x^2)=x+F^2$.
(3) Modulo~$2$ the equation reads $F+F^2F'=x$. In characteristic~$2$,
$F''=0$ because $n(n-1)$ is even. Differentiating gives $F'=1$, hence
$F=x+F^2$.
(4) Modulo~$2$, $(A-T_0)T_1=0$ with $T_1$ a unit, so $T_0=A$ and
$(A-1)A=x$. Then $B=A-1$ satisfies $B(0)=0$ and $B+B^2=x$, so $B=x+B^2$, and
$a(n)=b(n)$ for $n\ge1$.

For A273958, constant terms give $C(0)^2=0$. By Lemma~\ref{lem:dyadic},
$c=\overline C$ is the indicator of the powers of~$2$. Put $B=x\overline A$,
so that $b(n+1)=f(n)$ and $B+B^2=c^2$, i.e.\ $B(x)+B(x^2)=c(x^2)$. Hence
$b(n)=0$ for odd $n$ and $b(2m)=b(m)+c(m)$. The powers of~$2$ split into the
disjoint sets $\{4^j\}$ and $\{2\cdot4^k\}$, and strong induction gives
$b(n)=1$ exactly for $n=2\cdot4^k$.
\end{proof}

\section{Explicit solutions modulo a prime}
In the next results we exhibit the reduction $F$ explicitly and show that it
is the only solution by the $x$-adic observation of Section~2.

\begin{theorem}[A388734]\label{thm:388734}
If $A=1+xA^2+x^2(1-x)A^3$, then every $a(n)$ is odd.
\end{theorem}
\begin{proof}
Modulo~$2$, $U=1/(1-x)$ is a solution: $1+(x+x^2)/(1-x)^2=1/(1-x)$ because
$x+x^2\equiv x(1-x)$. Subtracting the equations for $F$ and $U$ gives
$F-U=(F-U)\,x\,[F+U+x(1-x)(F^2+FU+U^2)]$, so $x^n\mid F-U$ implies
$x^{n+1}\mid F-U$. Hence $F=U$.
\end{proof}

\begin{theorem}[A377100]\label{thm:377100}
If $a(0)=0$, $a(1)=1$ and $A=A(x^3)/A(x^2)+A^2$, then $a(n)\equiv1\pmod3$
for $n\ge1$.
\end{theorem}
\begin{proof}
Modulo~$3$ the equation reads $F(x^2)(F-F^2)=F(x^3)$. The series $E=x/(1-x)$
is a solution: both sides equal $x^3/(1-x)^3$, because $(1-x)^3=1-x^3$ and
$1-2x\equiv1+x$. For $D=F-E$ one computes
$E(x^2)\,D\,(1-F-E)=D(x^3)-D(x^2)(F-F^2)$. If $x^n\mid D$ with $n\ge2$, the
right-hand side is divisible by $x^{2n+1}$; since $E(x^2)$ is $x^2$ times a
unit and $1-F-E$ is a unit, $x^{2n-1}\mid D$. As $x^2\mid D$, we get $F=E$.
\end{proof}

\begin{theorem}[A274479]\label{thm:274479}
If $a(0)=0$, $a(1)=1$ and $A^2=A\bigl(x^2/(1-2x-4x^2)\bigr)$, then
$a(n)\equiv1\pmod3$ for $n\ge1$.
\end{theorem}
\begin{proof}
Modulo~$3$, $h=x^2/(1+x-x^2)$ and $E=x/(1-x)$ satisfies
$E(h)=x^2/(1+x-2x^2)=x^2/(1-x)^2=E^2$. For $D=F-E$ we have $D(F+E)=D(h)$. If
$x^n\mid D$ then $x^{2n}\mid D(h)$, and $F+E=xW$ with $W(0)=2$ a unit, so
$x^{2n-1}\mid D$. As $x^2\mid D$, we get $F=E$.
\end{proof}

\begin{theorem}[A120566]\label{thm:120566}
If $a(0)=0$, $a(1)=1$ and $A=A\circ A-x\,A\circ A\circ A$, then $a(n)$ is odd
for every $n\ge1$.
\end{theorem}
\begin{proof}
Modulo~$2$, $E=x/(1-x)$ satisfies $E\circ E=x/(1-2x)=x$, so
$E\circ E-xE\circ E\circ E=x-xE=E$. Let $D=F-E$ and suppose $x^n\mid D$ with
$n\ge2$. Since composition preserves congruences modulo $x^n$ and
$F=x+O(x^2)$, we have $D\circ F\equiv D$,
$E\circ F-E\circ E=D/((1-F)(1-E))\equiv D$, and
$x(F\circ F\circ F-E\circ E\circ E)\equiv0$, all modulo $x^{n+1}$. Since
$F\circ F-E\circ E=D\circ F+(E\circ F-E\circ E)$, subtracting the equations
for $F$ and $E$ gives $D\equiv2D\equiv0\pmod{x^{n+1}}$. As $x^2\mid D$, we get
$F=E$.
\end{proof}

\begin{theorem}[A361047]\label{thm:361047}
If $A(0)=0$ and $A(x-x^3A'^2)=x$, then $[x^m]A\equiv1\pmod3$ when $m$ is a
power of~$3$, and $[x^m]A\equiv0\pmod3$ otherwise.
\end{theorem}
\begin{proof}
Modulo~$3$ let $S=\sum_{k\ge0}x^{3^k}$. Then $S'=1$ and, by
Lemma~\ref{lem:frob}, $S-S^3=S-S(x^3)=x$, so $g\circ S=x$ for $g=x-x^3$.
Hence $S\circ g=x$, i.e.\ $S(x-x^3S'^2)=x$. For $D=F-S$ with $x^n\mid D$, the
inner series $g_F=x-x^3F'^2$ and $g_S=x-x^3S'^2$ differ by $x^3D'(F'+S')$,
which is divisible by $x^{n+2}$. Therefore
$0=F(g_F)-S(g_S)=D(g_F)+\bigl(S(g_F)-S(g_S)\bigr)\equiv D\pmod{x^{n+1}}$.
Hence $F=S$.
\end{proof}

\begin{remark}
The entry A361047 writes $A=\sum_{n\ge1}a(n)x^{2n-1}$ and conjectures
$a(n)\equiv1\pmod3$ exactly for $n=(3^k-1)/2$. With this indexing the
statement is contradicted by the data ($a(2)=1$). It is correct for the
index $i=n-1$, since $x^{2i+1}$ with $i=(3^k-1)/2$ is $x^{3^k}$.
Theorem~\ref{thm:361047} is this intended statement.
\end{remark}

\section{Additive series modulo three}
Call $G\in\F_3[[x]]$ \emph{additive} if $[x^m]G=0$ whenever $m$ is not a
power of~$3$. By Lemma~\ref{lem:frob}, $G^{3^j}=G(x^{3^j})$, so the
composition of two additive series is additive.

\begin{theorem}[A184894]\label{thm:184894}
Let $P_0=x$, $P_{m+1}=P_m+P_m^3$, and $a(m)=[x^{2m-1}]P_m$ for $m\ge1$.
\begin{enumerate}
\item If $2m-1$ is not a power of~$3$, then $3\mid a(m)$.
\item If $2m-1=3^j$, then $a(m)\equiv\binom mj\pmod3$.
\item $3\mid a(365)$, although $2\cdot365-1=3^6$.
\end{enumerate}
\end{theorem}
\begin{proof}
Modulo~$3$, $P_{m+1}=P_m+P_m(x^3)$, so the coefficients $c(m,k)=[x^k]P_m$
satisfy $c(m+1,k)=c(m,k)+[3\mid k]\,c(m,k/3)$ with $c(0,k)=[k=1]$. By
induction, $c(m,k)=0$ unless $k$ is a power of~$3$, and
$c(m,3^j)=\binom mj$. Finally $\binom{365}{6}\equiv0\pmod3$ by Lucas's
theorem, since $365=(111112)_3$ and $6=(20)_3$.
\end{proof}

The entry conjectures that $a(m)\equiv0\pmod3$ except at $m=(3^n+1)/2$.
Part~(1) proves this statement as written, namely vanishing outside the
exceptional indices. Part~(3) shows that the stronger reading, that $a(m)$ is
nonzero modulo~$3$ at every exceptional index, is false; the first failure is
at $n=6$. By Lucas's theorem, $a((3^n+1)/2)\not\equiv0$
exactly when every base-$3$ digit of $n$ except the last is at most~$1$.

\begin{theorem}[A107099]\label{thm:107099}
Let $A(0)=0$, $[x^1]A=1$ and $A\circ A=x+4x^3$. Then $3\mid[x^n]A$ whenever
$n$ is not a power of~$3$, and $3\mid[x^{729}]A$.
\end{theorem}
\begin{proof}
Modulo~$3$ the equation reads $F\circ F=x+x^3$. Suppose $m$ is not a power
of~$3$ and $f(j)=0$ for all non-powers $j<m$. Let $T$ be the truncation of
$F$ below degree $m$; it is additive, so $[x^m]T\circ T=0$. Since
$F=T+f(m)x^m+O(x^{m+1})$ and $f(1)=1$, we get $[x^m]F\circ F=2f(m)$, while
$[x^m](x+x^3)=0$; hence $f(m)=0$. Thus $F=\sum_kc_kx^{3^k}$ is additive, and
comparing coefficients of $x^{3^k}$ in $F\circ F=\sum_ic_iF(x^{3^i})$ gives
$c_0=1$ and $2c_k+\sum_{0<i<k}c_ic_{k-i}=[k=1]$ for $k\ge1$. This yields
$(c_1,\dots,c_6)=(2,1,1,2,1,0)$.
\end{proof}

The entry writes $A=\sum_{n\ge0}a(n)x^{2n+1}$ and conjectures that
$[x^n]A\equiv0\pmod3$ except at $n=3^k$. Theorem~\ref{thm:107099} proves this
statement as written. It also shows that the stronger reading, non-vanishing
at every power of~$3$, is false: $[x^{729}]A=a(364)\equiv0\pmod3$. The OEIS
b-file of A107099 already shows this zero.

\section{Verification}
\paragraph{Numerical checks.} For each entry, one of two checks was made.
Either a prefix of the sequence was recomputed from the defining equation and
compared with the b-file of the entry, or the b-file was substituted into the
defining equation and checked modulo two primes. The lengths and moduli vary
between entries and are recorded with the verification programs~\cite{repo}.
All checks agree with the theorems. They are cross-checks only; the theorems
rest on the proofs and on the Lean verification.

\paragraph{Formal verification.} All theorems are verified in Lean~4
(version 4.33.1, Mathlib~\cite{mathlib} commit \texttt{0df444a3}) for every
integer power series satisfying the stated hypotheses. The Lean files are
\texttt{Oeis389472Mod3} (Theorem~\ref{thm:389472}) and, for the other
entries, \texttt{HannaA}\emph{nnnnnn} with the A-number of the entry (the file
\texttt{HannaA338633} treats both continued fractions). A separate check confirms
that \texttt{A.subst h} denotes the composition $A(h)$. For example,
Theorem~\ref{thm:parity} for A273958 reads, in ASCII notation,
\begin{verbatim}
theorem hanna_a273958 (C A : PowerSeries Int) (hC : C = X + C ^ 2)
    (hA : X * A + X ^ 2 * A ^ 2 = C ^ 2) (n : Nat) :
    Odd (coeff n A) <-> exists k, n + 1 = 2 * 4 ^ k
\end{verbatim}
The project was rebuilt from a clean directory. The recorded axiom audits of
the main results (\texttt{\#print axioms}) report only the standard axioms
\texttt{propext}, \texttt{Classical.choice} and \texttt{Quot.sound}; no native
evaluation is used. The code and logs are available at~\cite{repo}.

\section*{Tool and computational resource disclosure}
This work was carried out between 25 and 26 September 2026 in a workflow
directed by the author, using two AI coding agents: OpenAI Codex (desktop
application; models \texttt{gpt-6-astra} and \texttt{gpt-6-sol}) and
Anthropic Claude Code (model Claude Opus~5.5). The author chose the research
programme and directed the agents. Codex found and formalised the proof of
Theorem~\ref{thm:389472}. Claude Code agents searched the OEIS for candidate
conjectures, found and formalised the other proofs, and drafted this text.
Each proof was checked by a second agent, and the agents searched the
literature for previous proofs.
The author is not a professional mathematician. The author has read the
paper and, using a side-by-side table, compared the statements of the Lean
theorems with the theorems stated here. The correctness of the results rests on the written proofs and on the Lean
formalisation; the numerical checks described in the previous section are
additional cross-checks. In addition,
each proof was checked by an AI agent other than the one that found it, and
the formal statements were compared with the OEIS entries by a second
agent. AI systems are not authors of this paper; the author takes full
responsibility for its content.

\begin{thebibliography}{9}
\bibitem{mathlib}
The mathlib Community, The Lean mathematical library, in \emph{Proceedings of
the 9th ACM SIGPLAN International Conference on Certified Programs and Proofs}
(CPP '20), ACM, 2020, \url{https://doi.org/10.1145/3372885.3373824}.
\bibitem{oeis}
OEIS Foundation Inc., \emph{The On-Line Encyclopedia of Integer Sequences},
\url{https://oeis.org}, entries A107099, A120566, A184894, A240998, A273958,
A274479, A295762, A301933, A338633, A338634, A361047, A377100, A388734 and
A389472 (P.~D.~Hanna), accessed 25--26 September 2026.
\bibitem{pf389472}
A.~Perez Fontelles, \emph{Every odd-indexed term of OEIS A389472 after the
first is even}, 31 August 2026, paper 1042 in~\cite{pfrepo},
\url{https://github.com/astrafala/Conjectures/blob/58cf4596af4541652fb191e177ef9b7771388d53/paper-sources/01001-01500/01042-PROOF.tex}.
\bibitem{pfrepo}
A.~Perez Fontelles, \emph{Proofs and disproofs of open OEIS conjectures},
repository \texttt{astrafala/Conjectures},
\url{https://github.com/astrafala/Conjectures}.
\bibitem{repo}
H.~Dong, \emph{ai4math-results}, repository,
\url{https://github.com/Horace-Maxwell/ai4math-results}; archived at Zenodo,
\url{https://doi.org/10.5281/zenodo.22970254}.
\bibitem{tru389472}
The Omega Institute, \emph{Hanna's Odd-Index Parity Conjecture},
\texttt{trureturing} repository, 2026,
\url{https://github.com/the-omega-institute/trureturing/blob/b5e234ba04d648411140fe79d6ec71b249bc30bb/Blueprint/D5/S1/Recurrence/Parity/SquareCubeSubstitutionOddIndexParity.md}.
\bibitem{trurepo}
The Omega Institute, \texttt{trureturing} repository,
\url{https://github.com/the-omega-institute/trureturing}.
\end{thebibliography}
\end{document}
